% ID: SEL-D0028-OB-E02
% ESTADO: validada. S029, cotejada con examen y criterios.
\textbf{a)} El límite izquierdo en $0$ es $1$ y el valor
y límite derecho son $-a$. La continuidad exige $\boxed{a=-1}$.
Las derivadas de los tramos son $4/(1-2x)^2$ y $2x-1$,
con pendientes laterales $4$ y $-1$ en el empalme.
Por tanto no existe derivada en $0$.

\textbf{b)} Con $a=2$, la función crece en $(-\infty,0)$,
decrece en $[0,1/2]$ y crece en $[1/2,\infty)$.
El cambio de signo de la derivada en $1/2$ da un mínimo
relativo en $\boxed{(1/2,-9/4)}$.
En $0$ hay un salto: los valores próximos a la izquierda
son mayores que $f(0)=-2$ y los de la derecha son menores.
Por eso $0$ no es un extremo. No hay máximos relativos.
